% Source: https://artofproblemsolving.com/community/c3h1692108_proof_that_positive_integers_dont_exist
% Stable identity: aops-topic-1692108-q01
\textbf{The correspondences are correct.} The map $m\mapsto\dfrac {m}{2}$ sends $2,4,6,\dots$ to $1,2,3,\dots$; every positive integer $k$ is the image of exactly one even number, namely $2k$. The map $m\mapsto m-1$ sends $2,4,6,\dots$ to $1,3,5,\dots$; every odd positive integer $j$ is the image of exactly one even number, $j+1$.

\textbf{The error.} The argument treats ``the number of elements'' of an infinite set as an ordinary number $N$ and then uses the rule of ordinary arithmetic that $N+N=N$ forces $N=0$. For finite sets this is correct, but these sets are infinite, and for infinite sets ``having the same number of elements'' only \emph{means} that there is a one-to-one correspondence between them. With that meaning, the fact that a set has the same size as a part of itself is not a contradiction; it is exactly what makes a set infinite. Here the positive integers really do have the same size as their proper subsets, the even numbers and the odd numbers, and also the same size as the union of the two, so in the arithmetic of infinite sizes $N+N=N$ holds without $N=0$. (A finite set can never be matched one-to-one with a proper part of itself, which is why the cancellation works only for finite sets.)

So the two correspondences are correct; the step ``$N+N=N$, hence $N=0$'' is the error, and positive integers do exist.
